\section{Simple types, adjunctions, and strict globular expansions}
\label{sec:I}
\providecommand{\sd}{\operatorname{sd}}
\providecommand{\Core}{\operatorname{core}}
\providecommand{\Pos}{\mathrm{Pos}}
\providecommand{\Sf}{\mathcal F}
Throughout this section a poset is regarded as a category with one arrow
$x\to y$ when $x\leq y$. Its order complex is denoted by $K(P)$, and
$\sd P$ is the poset of nonempty finite chains of $P$, ordered by inclusion.
A categorical homotopy is a finite zigzag of pointwise inequalities of
monotone maps. A beat point is an element whose strict lower set has a
greatest element or whose strict upper set has a least element; an empty
strict lower or upper set does not satisfy the relevant condition.

\begin{claim}{I.1}{Proved}{Adjunctions and finite cores}
For finite posets $P,Q$, the following conditions are equivalent: they are
connected by a zigzag of maps each having a left or a right adjoint; they
are categorically homotopy equivalent; their beat-point cores are
isomorphic. This includes the empty poset.
\end{claim}
\begin{proof}
Deleting a downbeat point $x$, with greatest strict predecessor $y$, has
retraction $r(x)=y$ and $r(z)=z$ otherwise. The inclusion $i:P-x\to P$
satisfies $i\dashv r$, since $ri=\id$ and $ir\leq\id$. For an upbeat
point the inequalities reverse, and $r\dashv i$. Repeated deletion
terminates at a poset with no beat points. All the comparison maps are
categorical homotopy equivalences.

If a finite beat-free poset $R$ has $h\leq\id_R$ and $h\neq\id_R$,
choose a minimal moved element $x$. Every $z<x$ is fixed, so
$z=h(z)\leq h(x)<x$; hence $h(x)$ is the greatest strict predecessor of
$x$, a contradiction. The dual argument excludes $h\geq\id_R$ with
$h\neq\id_R$. No nonconstant zigzag can therefore leave $\id_R$.
If two such posets $R,S$ are homotopy equivalent by $f,g$, the homotopies
$gf\simeq\id_R$ and $fg\simeq\id_S$ imply the equalities $gf=\id_R$
and $fg=\id_S$. Thus $f$ is an order isomorphism. Applying the core
retractions to any homotopy equivalence proves both uniqueness of the
core and the stated implication. Conversely isomorphic cores connect
$P$ and $Q$ by the adjunctions just constructed. Finally an adjunction
$f\dashv g$ always gives $\id\leq gf$ and $fg\leq\id$, so every
adjunction zigzag consists of homotopy equivalences. The empty case
follows because a nonempty poset has no map to the empty poset.
This supplies the particular form of the finite-space core theorem used
here; compare \cite[Section~2]{I_BarmakMinian2008}.
\end{proof}

\begin{claim}{I.1a}{Proved}{A finite collapse with a nontrivial core}
There is a specified 22-element beat-free poset $P$ whose nerve is
simplicially collapsible. Thus $P$ and a point have collapsible nerves
and nonisomorphic cores.
\end{claim}
\begin{proof}
Use six symbols $a,a',b,b',c,c'$. Let $K$ have the seven triangular
facets obtained by choosing one symbol from each pair, except $abc$.
Include every face of these facets. Define $P$ to consist of all six
vertices, all nine edges other than $ab,ac,bc$, and all seven facets,
ordered by inclusion. This describes every object, arrow and composite.
Each vertex is below at least two incomparable retained edges; every
retained edge has two vertices below and two distinct facets above; each
facet has at least two incomparable retained edges below. Consequently
no element is a beat point. The exact finite collapse certificate for
$K(P)$ is generated and independently checked at each removal by
\texttt{code/I\_finite\_checks.py}; the ordered pairs of deleted free
faces and their unique codimension-one cofaces are in
\texttt{code/outputs/I\_finite\_checks.json}. All pairs have precisely
these defining properties, and the remaining complex is one vertex.
The verification uses only the displayed subset order. In particular
it does not infer collapsibility merely from acyclicity. Since $P$ is
its own core, I.1 separates it from a point.
\end{proof}

\begin{claim}{I.2}{Refuted}{Adjunctions plus subdivision do not generate simple types}
The equivalence relation on finite posets generated by adjunctions and
last-element subdivision maps is strictly finer than simple homotopy
of nerves. Its exact description is
\[
 P\sim Q\quad\Longleftrightarrow\quad
 \Core(\sd^r P)\cong\Core(\sd^s Q)
 \quad\hbox{for some }r,s\geq0.
\]
The poset of I.1a and a point refute the proposed equivalence.
\end{claim}
\begin{proof}
The needed subdivision invariant is categorical contractibility:
\cite[Corollary~4.18]{I_BarmakMinian2012} states that a finite poset is
categorically contractible if and only if its subdivision is. We verify
the ingredients in the form used here. A beat deletion induces deletion
of a dominated vertex of the order complex; conversely a strong
simplicial collapse induces a categorical homotopy equivalence of face
posets. For finite complexes, a complex and its barycentric subdivision
are simultaneously strongly collapsible: take a minimal strong core
$L$, and observe during any proposed dominated-vertex collapse of
$\sd L$ that the vertices corresponding to vertices and maximal faces
of $L$ cannot be the first such vertices removed. A maximal face cannot
be dominated because all its vertices remain; a vertex cannot be
dominated because that would put another vertex in every maximal face
containing it, contrary to minimality of $L$. Thus the collapse can
reach a point only when $L$ is a point. Finally a beat-free poset has a
minimal order complex: if a vertex $x$ were dominated by $y$, then every
element comparable with $x$ would be comparable with $y$. If $x<y$,
choose a maximal $z$ in $[x,y)$. Every $w>z$ is comparable with $y$,
and maximality of $z$ forces $w\geq y$, making $z$ an upbeat point.
The case $y<x$ is dual. These are the arguments of Theorem~4.15 and
Proposition~4.17 of that source, and they prove the required invariant.
Adjunctions preserve categorical contractibility by I.1. I.1a therefore
cannot connect to a point by the proposed moves, although its nerve
collapses to a point.

For the asserted exact relation, subdivision preserves categorical
homotopy equivalences of finite posets. To see this directly, a
pointwise comparison of monotone maps can be performed one source
vertex at a time, in decreasing order for an upward comparison. The
induced simplicial maps of order complexes are then contiguous at each
step. A contiguity $u,v$ induces the pointwise comparisons
$\Sf(u)\leq[\sigma\mapsto u(\sigma)\cup v(\sigma)]\geq\Sf(v)$
on face posets. Thus subdivision preserves the equivalence. Any finite
zigzag can now be given common subdivision exponents: for example,
$\sd^a P\simeq\sd^b R$ and $\sd^c R\simeq\sd^d Q$ combine to
$\sd^{a+c}P\simeq\sd^{b+d}Q$. I.1 proves the core formulation.
Every allowed map is simple: an adjunction is a finite-poset homotopy
equivalence and hence a composite, up to homotopy, of beat deletions
and their inverses; subdivisions are simplicial subdivisions. This proves
strictness as well as failure of the converse.
\end{proof}

\begin{claim}{I.3}{Proved}{One-point weak equivalences generate simple types}
For finite posets $P,Q$, their nerves have the same simple homotopy type
if and only if they are connected by a zigzag of one-point inclusions
$R-x\hookrightarrow R$ that are weak homotopy equivalences.
Write $A=K(R-x)$ and
$L=K(R_{<x})*K(R_{>x})$, the link of $x$ in $K(R)$, with the usual
empty-complex convention. Such a one-point inclusion is a weak
equivalence exactly when $L$ is nonempty and integrally acyclic and,
in the component of $A$ containing $L$, the map
$\pi_1 L\to\pi_1 A$ is trivial. Every such inclusion has zero
Whitehead torsion. Among these inclusions, the sieve Dwyer maps are
exactly deletion of a maximal downbeat point. Dually the cosieve
version is deletion of a minimal upbeat point.
\end{claim}
\begin{proof}
We have $K(R)=A\cup_L(x*L)$. If $A\hookrightarrow K(R)$ is a weak
equivalence, the relative integral homology groups vanish, and
$H_j(K(R),A)\cong\widetilde H_{j-1}(L)$ shows that $L$ is nonempty,
connected and acyclic. The image of its fundamental group is killed
by its cone, so injectivity of $\pi_1A\to\pi_1K(R)$ forces that
image to be trivial. Conversely these conditions and van Kampen give
an isomorphism on fundamental groups and on components. Over the
universal cover of the affected component, the lifted relative chain
complex is
\[
 \mathbb Z[\pi_1A]\otimes_{\mathbb Z}
 \widetilde C_{*-1}(L;\mathbb Z).
\]
It is acyclic, so the simply connected homology criterion proves a
weak equivalence. The same formula computes torsion. An acyclic finite
free integral chain complex has torsion represented by integral
unimodular matrices, whose classes vanish after passage from
$K_1(\mathbb Z)$ to $\Wh(\pi_1A)$. Extension of scalars therefore gives
zero torsion. Other components are untouched. This also explains why
one must not replace the condition by contractibility of the link:
acyclic non-simply-connected links are permitted.

The forward direction for a zigzag follows by composing these simple
equivalences and their inverses. For the reverse direction use
\cite[Theorem~3.10(a)]{I_BarmakMinian2008}, whose exact finite-poset
statement is that simple homotopy of order complexes is generated by
adding and removing weak points, meaning points with a categorically
contractible strict upper or strict lower poset. These are among the
inclusions in the statement: their links are joins with a contractible
factor. The theorem's construction applies to every finite poset;
it passes from a simplicial elementary collapse to its face posets
by deleting the free face and then its unique maximal coface, and
connects a poset to its face-poset subdivision by weak-point moves.
No assumption that $P$ or $Q$ themselves are face posets is imposed.

For the last assertion, $R-x\hookrightarrow R$ is a sieve precisely
when $x$ is maximal. The nonempty link implies that $x$ has an old
predecessor, so the upward closure of $R-x$ is all of $R$. A right
adjoint retraction to the inclusion then exists exactly when the
strict lower set of $x$ has a greatest element. This is the definition
of a downbeat point. If the deleted point is isolated the inclusion
can be Dwyer using a smaller cosieve, but it is not a weak equivalence;
thus it does not supply an omitted case. The dual statement follows
by reversing the order.
\end{proof}

\subsection{Realization and the precise invariance of torsion}
We use the convention in which a relative cellular complex concentrated
in degrees $3,2$, with differential an invertible matrix $M$, represents
the torsion $[M]\in\Wh(\pi)$. Reversing the global torsion convention
changes the signs in the realization assertion, without changing any
zero or nonzero assertion below. Fundamental groups are based and
identified with the indicated group; changing base paths acts by inner
automorphisms, which act trivially on the Whitehead group.
The chain definition and its additivity are those of
\cite[Definitions~12 and~17, Exercise~14]{I_LurieWhiteheadI};
the equivalence between zero torsion and simple homotopy equivalence
is \cite[Corollary~5 and Theorem~8]{I_LurieWhiteheadII}.

\begin{claim}{I.4}{Proved}{Every Whitehead class is realized by finite posets}
For every finitely presented group $\pi$ and every
$\tau\in\Wh(\pi)$, there is a weak equivalence of connected finite
posets $f:P\to Q$, with specified identifications
$\pi_1K(P)\cong\pi\cong\pi_1K(Q)$, whose Whitehead torsion is $\tau$.
Every composite of the maps in I.2 has zero torsion. More generally,
every zigzag of the moves in I.2, interpreted as a homotopy equivalence
of marked finite complexes, has zero torsion.
\end{claim}
\begin{proof}
Choose a finite presentation complex $B$ for $\pi$ and an invertible
$r\times r$ matrix $M$ over $\mathbb Z\pi$ representing $\tau$.
The empty matrix deals with $\tau=0$. Attach $r$ two-spheres at the
basepoint, giving $B'=B\vee\bigvee_{j=1}^r S^2_j$.
In the universal cover the new spheres freely generate the summand
$(\mathbb Z\pi)^r$ of $\pi_2(B')$. For each column of $M$, represent
the corresponding element by a map $\alpha_j:S^2\to B'$:
pinch a finite collection of disjoint disks to spheres, give each
sphere the required integer degree, and join its basepoint to the
global basepoint along a loop representing the required group element.
All coefficients and loops are finite. Attach a three-cell along each
$\alpha_j$, obtaining the finite CW pair $(Y,B)$.

The inclusion $B\hookrightarrow Y$ induces an isomorphism on
fundamental groups, since the new cells have dimensions two and three
and the two-cells were attached by constant maps. Its relative
universal-cover chain complex is exactly
\[
 0\longrightarrow (\mathbb Z\pi)^r
   \xrightarrow{\ M\ }(\mathbb Z\pi)^r\longrightarrow0,
 \qquad\text{in degrees }3,2.
\]
Invertibility of $M$ makes this acyclic. The induced map of simply
connected universal covers is a homology equivalence, hence a homotopy
equivalence by \cite[Corollary~4.33]{I_Hatcher2002}. The induced
map on fundamental groups is already an isomorphism, so
\cite[Theorem~4.5]{I_Hatcher2002} applies downstairs.
Thus $B\hookrightarrow Y$ is a homotopy equivalence and has torsion
$[M]=\tau$ by the definition of cellular torsion.
This is the matrix realization of
\cite[Remark~6]{I_LurieWhiteheadII}, with its dimensions and
marking specified here.

We specify why this gives a finite simplicial, and then a finite-poset,
realization without losing the torsion. Replace the presentation
circles and disks by finite polygonal models, and replace every
attaching map by a simplicial approximation on a finite subdivision of
its source \cite[Theorem~2C.1]{I_Hatcher2002}.
Build attachments by mapping cylinders followed by cones.
A finite simplicial map has a finite polyhedral mapping cylinder, and
the mapping-cylinder replacement of an attaching map collapses onto
its original target. Replacing a map by a homotopic map changes its
mapping cylinder by a simple equivalence relative to the target. Both
assertions follow by triangulating the prism of the homotopy and
removing its simplices successively from a free end. Consequently we
obtain a finite simplicial pair $(Y_s,B_s)$ with simple identifications
to $(Y,B)$ and the same torsion.
Equivalently, for two homotopic attaching maps, attach their two
disks and the intervening prism as one cell of the next dimension;
either disk is a free face of that new cell. The two resulting
inclusions are elementary expansions relative to the old target.
This is the explicit comparison in
\cite[Example~11]{I_LurieWhiteheadII}, where the homotopy between
the attaching maps is required.

For completeness, the same construction can be made entirely with
finite orders. For a monotone map $u:S\to T$ let $B(u)$ be the disjoint
union $S\amalg T$ with their old orders and new comparisons
$s\leq t$ exactly when $u(s)\leq t$. Deleting the elements of $S$
in reverse linear order contracts $B(u)$ onto $T$ by upbeat deletions.
The comparison from the ordinary mapping cylinder of $K(u)$ to
$K(B(u))$, supplied by $s\leq u(s)$, is the identity at both ends.
Its torsion is zero: both inclusions of $K(T)$ are simple, so the
composition formula for torsion gives this assertion. Coning off the
source preserves this simple comparison. Indeed the relative cellular
chain matrices of the glued map are block triangular, with the old
comparison on one diagonal block and the identity on the cone block;
the torsions therefore add, with the overlap subtracted, and all three
terms are zero. This also proves the assertion after the fundamental
group changes under coning, by extension of scalars along the induced
group-ring map. Applying this recipe to face-poset maps is a finite
algorithm for the simplicial replacement just used.

Finally take $P=\Sf(B_s)$, $Q=\Sf(Y_s)$ and let $f$ be the inclusion
of face posets. Their order complexes are barycentric subdivisions,
and the canonical subdivision homeomorphisms are simple: stellar
subdivision of a face is obtained by an elementary expansion followed
by elementary collapses, and barycentric subdivision is the finite
iteration in decreasing face dimension. Hence $\tau(f)=\tau$.

A categorical homotopy equivalence of finite posets has zero torsion:
the core retractions of I.1 are beat deletions, their order complexes
collapse, and the induced map between cores is homotopic to an
isomorphism. A subdivision map also has zero torsion by the preceding
subdivision argument. The formulas
\[
 \tau(gf)=\tau(g)+g_*\tau(f),\qquad
 \tau(f^{-1})=-f_*^{-1}\tau(f)
\]
now prove the last two assertions, component by component if needed.
\end{proof}

\begin{claim}{I.5}{Refuted as literally worded}{Torsion is an invariant of marked homotopy classes}
Whitehead torsion of a weak equivalence between specified finite posets
is invariant under natural transformations and under localization at
the one-point maps of I.3. It is also invariant under equality of its
fixed-source, fixed-target class in $\Cat[\Winfty^{-1}]$.
Thus the additional assertion ``and on no fundamental localizer'' is
false if it denies this last invariance. The correct obstruction is
different: the class of simple equivalences of finite posets cannot be
the restriction of a fundamental localizer contained in $\Winfty$.
No fundamental localizer refines $\Winfty$ so as to distinguish simple
equivalences from other weak equivalences.
\end{claim}
\begin{proof}
A natural transformation between poset maps realizes to a homotopy.
Cellular torsion of homotopy equivalences is invariant under homotopy:
the mapping cones of the induced chain maps are related by a
block-unipotent chain isomorphism, of torsion zero, using the chain
homotopy. Thus the assertion includes a zigzag of transformations.
This is also \cite[Proposition~1]{I_LurieWhiteheadII}; its
Proposition~4 gives the composition formula in I.4, after transporting
the fundamental-group marking.
Realization sends every inverted map of I.3 to a homotopy equivalence
of torsion zero, so it factors through the localization as a functor
to the homotopy category of spaces. Equality of two arrows in this
localization implies equality of their realized homotopy classes;
their torsions are therefore equal under the specified markings.
For a zigzag, torsion can equivalently be computed by the composition
and inverse formulas of I.4; an inverted elementary map contributes
zero. These statements require neither a calculus of fractions nor a
strict representative of each localized arrow.

The identical reasoning works for equality of two arrows in
$\Cat[\Winfty^{-1}]$: this category is equivalent, by nerve and
realization, to the homotopy category of spaces, and maps between the
fixed finite CW complexes $K(P),K(Q)$ have their ordinary homotopy
classes there. This supplies a fundamental localizer for which the
literal negative assertion fails.

On the other hand, the minimal-fundamental-localizer theorem gives
$\Winfty\subseteq W$ for every fundamental localizer $W$ in the
relative-Theorem-A sense \cite[Corollary~4.2.19]{I_Cisinski2006}.
If also $W\subseteq\Winfty$, equality is
forced. I.4 supplies finite-poset weak equivalences with nonzero
torsion whenever $\Wh(\pi)\ne0$; the explicit $C_5$ example in I.8b
below avoids making that nonvanishing conditional. Such a map must
belong to $W$, so the simple equivalences cannot be its restriction.
This confirms the exact obstruction in POLES D.4--D.6: a class of
maps declared invertible by a fundamental localizer cannot encode the
extra finite simple structure. It does not say that torsion of a
fixed marked map ceases to be homotopy invariant.
\end{proof}

\subsection{What the strict globular moves actually preserve}
The disks in 312 are free disks of a Grothendieck coherator model,
whereas those used below are free \emph{strict} globes. There is no
identification of these two categories of models. We make the strict
interpretation explicit because it is the one in the proposed
pushouts in $\Cat$ and in strict $\omega$-categories.

\begin{claim}{I.312}{Proved}{Audit of the coherator expansion and its comparison theorem}
For the Grothendieck coherators defined in the supplied 312 source,
the elementary expansion of any cellular model along an arbitrary
cell is a weak equivalence. Its left semi-model consequence and
the deduction of the homotopy hypothesis from Henry's comparison
are valid with the stated cofibrant-source restrictions. These
statements concern weak coherator groupoids and do not identify
their disks or pushouts with those of strict categories.
\end{claim}
\begin{proof}
We checked the four proof sections of 312 and both attached course
chapters. Here are the separate steps needed to confirm the
argument. In its preliminaries, division is performed on homotopy
classes of cells with \emph{specified} parallel boundaries. The
successive boundary corrections are operations on globular pasting
schemes, filled by admissible pairs of the coherator. They provide
left inverses for multiplication; all the multiplication maps are
injective by the same correction, so these left inverses also give
the required right inverses. This remains valid when a relevant
hom-set is empty. Whiskering to iterated units gives the
fixed-boundary form of the weak-equivalence test. Surjectivity is
the existence of a cell with the prescribed boundary up to one
higher homotopy, and injectivity is the same test in the next
dimension. Thus the test is equivalent to bijection on components
and isomorphisms of based homotopy groups. Its witnesses use
finitely many cells. This proves the filtered-colimit and retract
closures used later. The lifting calculation gives
$I\text{-}\mathrm{inj}=W\cap J\text{-}\mathrm{inj}$: lift a disk
up to homotopy by the test, then lift that homotopy along $J$ to
correct the disk exactly. These are the respective claims of
312, Section~2 (source file \texttt{01-preliminaries.tex}).

The cutoff argument in Section~3 retains cells through dimension
$q$ and gives every parallel higher boundary exactly one filler.
Its interpretation of operations follows the \emph{original}
free order of the coherator generators. If an operation has output
$j\leq q$, its generating arity has dimension at most $j$, so
its tuple is an unchanged tuple of the original model. Its faces,
even if represented by composites with a larger apparent arity,
are already respected by induction on that original order. Above
$q$, uniqueness of the filler supplies both the operation and
compatibility. This proves a genuine reflector preserving all
cells through $q$; it is not a reordering of generators by
dimension. As a left adjoint it preserves the algebraic pushouts
and cellular presentations under discussion.

At a fixed cutoff the proof descends in the expansion dimension.
The top-dimensional assertion holds because the new parallel
$q$-cell is already joined to the old one by the unique higher
filler, while lower cells do not change. In the upper-cell lemma,
an upper boundary-cell weak equivalence is factored by upper
expansions followed by their right lifting class. The descending
induction makes the first factor a weak equivalence. The exact
boundary test makes the second factor lift the upper boundaries,
so the first map is a retract of an upper expansion-cell map.
Its splitting and pushout assertion follow from that retract.
The free-shape assertion uses its finite low-dimensional spine;
higher attachments leave every fixed lower-dimensional test
unchanged. For tree-shaped diagrams, the contraction removes a
last branch and the backtracking comparison is supplied in
dimension two higher. Each use of an expansion in these arguments
has dimension strictly above the current induction index. Thus
these shape and tree assertions do not assume the expansion being
proved at that index.

Section~4 constructs a partial cylinder with $P_j=D_j$ below the
current index $n$, $P_n=D_{n+1}$, and upper boundary-cell completions
thereafter. The skeletal gluing formula makes each one-ended
inclusion an upper boundary-cell weak equivalence, hence split by
the preceding lemma. Its realization on every generating pasting
scheme is contractible by the tree assertion. The path object is
then interpreted operation by operation in the original coherator
order. For an operation of output $j$, its two endpoint values and
already defined boundary operations give a map from the
two-ended boundary into that contractible realization. The disk
lifting property supplies its extension. This verifies the full
operations and their prescribed faces, not just a map of the
underlying globular sets. The endpoint maps lift all boundaries.
For the joint endpoint map, the displayed tree domain in the
source is an upper boundary-cell weak equivalence into $P_{n+1}$;
its splitting supplies the required lift against $J_n$.

We also checked the final arrow-retract argument, where a mistaken
use of cancellation would be consequential. For
$i:X\to Y=X\amalg_{D_n}D_{n+1}$ the unit gives a strict
retraction $r:Y\to X$. Choose a section $h:X\to PX$ of the
first endpoint and put $u=p_1h$. Both endpoints are trivial
fibrations, so $u\in W$. The joint-endpoint lift gives
$v=iur\in W$. The arrow $iu:X\to Y$ is a retract of
$v:Y\to Y$ by domain maps $i,r$ and identity codomain maps:
\[
 vi=iu,\qquad (iu)r=v,\qquad ri=\id_X.
\]
Thus $iu\in W$ and then $i\in W$. For a test in dimension $k$
before cutoff, choose $q\geq\max(n+1,k+1)$; its disk and its
one-higher homotopy are in unchanged dimensions. The cutoff result
therefore proves exactly the original test, for an arbitrary
attaching cell, including a composite or a coherence cell.

The semi-model transfer of Section~5 uses no expansion with a
noncofibrant source. A cofibrant model is an arrow retract of a
cellular one. An expansion preserves that retract by using the
identity on its newly attached disk. A relative $J$-cell map with
cofibrant source is consequently a weak equivalence, at successors
by the expansion theorem and at limits by finite-witness closure.
The small-object argument gives the two factorizations, with the
second required only for cofibrant sources. The identity
$I\text{-}\mathrm{inj}=W\cap J\text{-}\mathrm{inj}$ and the
arrow-retract argument give the corresponding lifting axioms.
Units fill every $J$-lifting problem to the terminal object, so
all objects are fibrant. This is precisely the left semi-model
statement asserted by the course.

The external comparison is explicit. Henry's Conjecture~5.3.3 is
the finite cellular expansion assertion just obtained. Under that
hypothesis, Proposition~5.3.6 makes the finite cellular models a
cylinder category; Proposition~5.3.7 identifies its completion with
the models of the same coherator; Lemmas~5.3.9 and 5.3.11 give
the lifting against fibrations and the homotopy-terminal
condition; Corollary~5.3.12 makes it a cylinder coherator.
Theorem~5.1.2(5) then states that its completion has a
geometric-realization/fundamental-groupoid Quillen equivalence
with spaces. Its proof uses the explicit semisimplicial cylinder
coherator and Theorem~5.2.10, whose realization is acyclic:
homotopy full faithfulness is Proposition~5.2.8 and essential
surjectivity follows by finite cell induction. These exact
statements and proofs were opened in the primary source
\cite{I_Henry2016}. Corollary~5.3.13 records their implication for
Grothendieck groupoids. The finitely cellular hypothesis is
satisfied by the all-cellular theorem of 312. We use this
comparison with its hypothesis discharged, not the homotopy
hypothesis as a premise. In Lemma~5.3.11 of the opened version,
the map from the homotopy slice is called a ``trivial cofibration'';
Proposition~4.4.4 identifies the needed map as an \emph{acyclic
fibration}, and the lifting argument uses the latter property.
This is a terminological correction in that paragraph, not a
missing comparison assumption.
\end{proof}

\begin{claim}{I.6}{Refuted}{The four-point diamond obstructs the proposed moves in Cat}
Let $\tau_{\leq1}:\omega\Cat\to\Cat$ be the left adjoint to inclusion,
and put $D_j^{(1)}=\tau_{\leq1}D_j$. Pushouts in $\Cat$ of
$D_0^{(1)}\to D_1^{(1)}$ and $D_1^{(1)}\to D_2^{(1)}$, and inverse
moves, do not generate all adjunction zigzags between finite posets.
This remains false if both source and target face inclusions are
allowed. The four-point diamond $0<1,2<3$ is a counterexample.
\end{claim}
\begin{proof}
Reflection forces every 2-cell to be an identity, identifying its
parallel source and target arrows. Hence $D_j^{(1)}\cong[1]$ for
$j\geq1$, and the reflected face maps between them are identities.
A source attachment $D_0^{(1)}\to D_1^{(1)}$ at $x\in P$ adds a
fresh maximal point $y$ with strict lower set $\mathord\downarrow x$.
It adds precisely the Hasse edge $x<y$ and changes no old Hasse edge.
The target version dually adds a fresh minimal leaf. Thus all
nontrivial moves are additions or removals of leaves of the oriented
Hasse graph. For a finite poset put
\[
 \lambda(P)=|\operatorname{Cov}(P)|-|P|+|\pi_0K(P)|.
\]
This is the cycle rank of its undirected Hasse graph and is unchanged
by leaf moves. The diamond has four vertices, four cover relations,
and one component, so $\lambda=1$; a point has $\lambda=0$.
Nevertheless deleting the internal beat point $1$, then the resulting
chain, shows that the diamond has point core. Its nerve is a cone
because it has a least element. The finite order, core and cycle-rank
calculations are checked in \texttt{I\_finite\_checks.py}.

With both endpoints allowed, the exact relation produced by these
moves is equality of the reduced oriented Hasse graphs obtained by
successively pruning leaves, retaining a single vertex in a tree
component. This reduced graph is independent of pruning order:
vertices on cycles and the paths between cycles survive; every other
edge lies in a tree attached to this subgraph, or in a tree component,
and can be removed from its ends. Isomorphic reduced diagrams give
back the original posets by the inverse leaf additions. With source
faces only the relation is still smaller, permitting only maximal
leaf additions and deletions. Neither variant is the core relation
of I.1.
\end{proof}

\begin{claim}{I.7}{Refuted}{Higher strict globular expansions do not repair the obstruction}
For finite posets $P,Q$, simple homotopy of $K(P),K(Q)$ does not imply
that $c_\infty NP,c_\infty NQ$ are connected by pushouts of strict
$D_n\to D_{n+1}$ and their inverse moves, even when all dimensions
and both face orientations are permitted. Again the diamond and a
point give a counterexample.
\end{claim}
\begin{proof}
The reflector $\tau_{\leq1}$ preserves pushouts. For every ordinary
category $P$,
\[
 \tau_{\leq1}c_\infty NP=c_1NP\cong P.
\]
For $n\geq1$ it sends $D_n\to D_{n+1}$ to the identity of $[1]$.
For $n=0$ its effect is the leaf attachment just computed. Therefore
any purported zigzag would yield a zigzag in $\Cat$ whose only
nonisomorphism steps add or remove a Hasse leaf. Starting at a poset,
all terms remain posets: a leaf attachment preserves being a poset,
and its source is a full subcategory of its target, so the same holds
when the move is traversed backwards. The invariant $\lambda$ of
I.6 excludes the diamond-to-point zigzag. Both starting nerves are
simplicially collapsible, so every required hypothesis is verified.
This obstruction concerns existence of the zigzag and hence is
independent of the model-structure assertion of 317. The strongest
uniform necessary condition supplied by this argument is the precise
leaf-equivalence condition on the reflected 1-categories; mere simple
homotopy type does not imply it. General finite polygraphic elementary
expansions must allow a free face in a composite boundary, not just
one entire globular source face, if they are to avoid this obstruction.
\end{proof}

\subsection{Cofibrations and an explicit torsion certificate}
\begin{claim}{I.8a}{Refuted}{Dwyer maps strictly contain the finite-poset cofibrations}
Every Thomason cofibration between posets is a Dwyer map. The converse
fails, including for finite posets and for retracts of Dwyer maps.
Let $T$ have six elements $a_1,a_2,b_1,b_2,c_1,c_2$, with
$a_i<b_j<c_k$ for every $i,j,k$ and no other strict comparisons except
those forced by transitivity. Then $\varnothing\hookrightarrow T$ is
Dwyer but is not a Thomason cofibration.
\end{claim}
\begin{proof}
Generating Thomason boundary maps are Dwyer maps; their pushouts and
transfinite composites are Dwyer maps. For posets, retracts remain
Dwyer. To check the latter point explicitly, write an arrow retract
$i:A\to B$ of $j:C\to D$, with sections $A\to C$, $B\to D$ and
retractions in the other direction. Fullness and the sieve condition
pass to the retract. Let $Z\subset B$ be the upward closure of $A$.
The section takes $Z$ into the upward closure of $C$, where the Dwyer
right adjoint is defined. Compose the section, this right adjoint,
and the retraction $C\to A$. The resulting monotone map $r:Z\to A$
satisfies $ri=\id$ and $ir\leq\id$. These identities give $i\dashv r$.
Thus every cofibration of posets is Dwyer, as in
\cite[Lemma~3.13]{I_MayStephanZakharevich2017}.

Every initial map $\varnothing\to T$ is Dwyer: its cosieve
neighborhood may be empty. We give the noncofibrancy proof, rather
than relying only on the exhibited six-element diagram in
\cite[Proposition~6.2]{I_MayStephanZakharevich2017}.
Consider attachment of one generating boundary cell
\[
 Y=X\amalg_{c\sd^2\partial\Delta^n}c\sd^2\Delta^n.
\]
An object of $c\sd^2\Delta^n$ is a nonempty flag of nonempty subsets
of $\{0,\ldots,n\}$, and comparison means inclusion of flags.
The new objects are exactly flags containing the full subset $V$.
Apart from the singleton flag $z=\{V\}$, deleting $V$ defines an
order-preserving map
$p:Y-\{z\}\to X$ which is the identity on $X$ and satisfies
$p(y)\leq y$. Mixed comparisons are exactly those obtained by first
passing to this old flag.

Suppose $T$ is an order retract of $Y$, via $i:T\to Y$, $r:Y\to T$.
If both $i(b_1),i(b_2)$ were new, their union, as flags, would be a
new flag contained in both $i(c_1),i(c_2)$. Applying $r$ would give
an element above both $b_j$ and below both $c_k$, which does not exist
in $T$. Thus some $i(b_j)$ is old. If either $i(a_j)$ were new, all
four images of the $b$'s and $c$'s would be new, giving the same
contradiction. Hence both $i(a_j)$ are old. In particular $z$ is not
in $i(T)$: it is a minimal new object with no old predecessor, so only
a minimal element of $T$ could map to it.

The endomorphism $rpi:T\to T$ is defined and is below the identity.
It fixes $a_1,a_2$ and an old $b_j$. For the other $b_j$ its value is
above both $a$'s and below $b_j$, forcing that value to be $b_j$.
Its values at the $c$'s are then above both $b$'s and below the
corresponding $c$, so these are fixed too. Therefore $rpi=\id_T$:
any retraction of $T$ from $Y$ already gives one from $X$.

If $T$ were cofibrant, it would be a retract of a transfinite
composition of such cell attachments starting at the empty poset.
Choose the first stage admitting the finite image of the section
$T\to Y$. It is not a nonzero limit stage, since that finite image
is contained in an earlier stage. The preceding argument excludes a
successor stage. The empty stage cannot contain $T$. This
contradiction proves noncofibrancy. The displayed order is explicitly
checked by \texttt{I\_finite\_checks.py}; the proof uses no unstated
choice of arrows in the six-element category.
\end{proof}

\begin{claim}{I.8b}{Refuted}{Finite Thomason trivial cofibrations can have nonzero torsion}
Every class $\tau\in\Wh(\pi)$ for a finitely presented group $\pi$
is realized by a Thomason trivial cofibration between finite posets.
In particular there is an explicitly specified example with
$\pi=C_5$ and nonzero torsion. The program
\texttt{code/I\_torsion\_witness.py} and its JSON output give all
finite simplicial and order data for this example.
\end{claim}
\begin{proof}
Take the finite simplicial inclusion $B_s\hookrightarrow Y_s$
constructed in I.4. It is an acyclic cofibration in simplicial sets.
The left Quillen functor $c\sd^2$ therefore takes it to a Thomason
trivial cofibration. Both resulting categories are finite posets,
whose nerves identify with the second subdivisions of the original
complexes. Subdivision changes none of the torsion under the canonical
markings, so the result realizes $\tau$. This uses only the
dimension-one case of 317, namely the model structure and generating
maps of \cite{I_Thomason1980}.

Here are finite data and independent checks for a nonzero instance.
Let $C_m$ denote the simplicial polygon with vertices
$0,\ldots,m-1$ and edges $\{i,i+1\}$, indices modulo $m$.
The simplicial map $C_{15}\to C_3$ taking $i$ to $i\bmod3$ has
degree five. Its face-poset map has 30 source and 6 target objects.
Form its order mapping cylinder and add a new element above every
source object, and incomparable with all target objects. Call the
resulting 37-element poset $P_0$. By the mapping-cylinder comparison
in I.4, its order complex $B$ is a finite model, with its simple
structure, of the presentation complex $\langle t\mid t^5\rangle$.
In particular $\pi_1B=C_5$.

Wedge a tetrahedral boundary $S^2$ onto $B$ at its target-circle
vertex $0$. The program gives a simplicial map
$\alpha:S\to B\vee S^2$, with $S$ a specified triangulated
two-sphere. It checks that every edge lies in precisely two
triangles, that all vertex links are connected cycles, that the
triangles admit a consistent orientation, and that the connected
surface has Euler characteristic two. The surface classification
therefore identifies it as $S^2$; no homology-sphere recognition is
being assumed in higher dimension. Its numbers of vertices, edges
and triangles are respectively $388,1158,772$.

The attaching map is formed from nine degree-one or degree-minus-one
sphere pieces, connected by triangulated tubes following the powers
of the target-circle loop. The output records all vertex images and
all oriented triangles. It also records a $\mathbb Z/5$ edge cocycle
of $B$, of value one on $t$, and the induced deck label at every
source vertex. Counting the lifted images of one oriented face of the
target sphere gives the relative cellular boundary
\[
 v=2-t-3t^2-t^3+2t^4\quad\text{in }\mathbb Z[C_5].
\]
This is verified directly from the finite map, independently of the
intended tube construction. The multiplication identity
\[
 v(-1-t+2t^2-3t^3+2t^4)=1\pmod {t^5-1}
\]
is checked exactly over the integers. Thus attachment of the
three-cell has an acyclic relative universal-cover complex.
The augmentation of $v$ is $-1$, as is the determinant of its
regular-representation matrix; these provide additional checks.

Let $\zeta=e^{2\pi i/5}$. In $\mathbb C$ one has
$v(\zeta)=-(1+\zeta)^4$, because the difference is twice
$1+\zeta+\zeta^2+\zeta^3+\zeta^4$. Consequently
\[
 |v(\zeta)|=(2\cos(\pi/5))^4>1.
\]
The determinant after the ring homomorphism
$\mathbb Z[C_5]\to\mathbb C$, $t\mapsto\zeta$, kills elementary
matrices. Every trivial Whitehead unit $\pm t^j$ has complex modulus
one. Hence $[v]\ne0$ in $\Wh(C_5)$.

To give the actual finite posets, put $X=\Sf(B\vee S^2)$ and let
$s=\Sf(\alpha):\Sf(S)\to X$. Define $Y=B(s)\cup\{\ast\}$ by
placing every source element below $\ast$ and leaving $\ast$
incomparable with $X$. There are $266+2318+1=2585$ elements.
The map $A=\Sf(B)\hookrightarrow Y$ has, by the simple
mapping-cylinder comparison, the torsion $[v]$ just computed.
The desired trivial cofibration is
\[
 \sd^2 A\hookrightarrow\sd^2Y.
\]
An element of either of these two finite posets is explicitly a
nonempty inclusion-chain of nonempty chains in the indicated base
poset. Thus the JSON specification gives their complete object and
order sets without printing a redundant list of all double flags.
The verification enumerates the base complexes and the attaching map;
the last subdivision step follows from the exact displayed formula
and the left Quillen property. It is not a claim that all double flags
were materialized in the computation.
\end{proof}

\subsection{Finite fibrancy and the countable rigidity question}
\begin{claim}{I.9a}{Proved}{Every finite fibrant component has point core}
Every connected finite Thomason-fibrant poset is categorically
contractible. In particular a finite Thomason-fibrant poset with weakly
contractible nerve has a point as its core. More strongly, a connected
finite Thomason-fibrant poset with no upbeat points has a least element.
The empty poset is allowed as a disjoint union of zero such components.
\end{claim}
\begin{proof}
We first construct a finite test diagram. Let $P$ be a nonempty finite
connected poset with no upbeat points. Take one vertex $v_x$ for each
$x\in P$ and join these vertices by a spanning tree of the undirected
comparability graph of $P$, orienting each selected edge according
to the order in $P$. Label $v_x$ by $x$. At each vertex labelled $x$,
attach a fresh upward tree having a child labelled $z$ for every
$z>x$, and repeat this operation at each fresh child. Strict increase
of labels makes the recursion terminate. Only fresh vertices are used,
so the resulting undirected graph is a finite tree. Orient every new
edge from its parent to its child, and let $T$ be the order generated
by these arrows. Its label map $a:T\to P$ is monotone and onto.
Every vertex $t$ labelled $x$ has, for every $z>x$, a specified
successor labelled $z$. The poset $T$ is categorically contractible:
a leaf is a beat point, and deleting leaves reduces it to a point.

Put $D=\sd^2T$ and write
$\mu:D\to T$ for the twice-iterated last-element map. We claim that
\[
 h=a\mu:D\to P,\qquad g\geq h\quad\Longrightarrow\quad g=h.
 \tag{I.9.1}
\]
An element $\sigma\in D$ is a nonempty flag of nonempty chains in
$T$. Let $C$ be its largest chain, let $t=\max C$, and put
$x=a(t)=h(\sigma)$. We argue downward over the finite order $P$.
If $x$ is maximal, $g(\sigma)\geq x$ implies equality.
For every $z>x$ choose the successor $t_z>t$ labelled $z$ and form
\[
 \sigma_z=\sigma\cup\{C\cup\{t_z\}\}\in D.
\]
It is a larger flag, and $h(\sigma_z)=z$. The induction hypothesis
gives $g(\sigma_z)=z$, so
$x\leq g(\sigma)\leq z$ for every $z>x$.
If $g(\sigma)>x$, it would therefore be the least strict upper
element of $x$, contrary to the no-upbeat assumption. This proves
(I.9.1) for every flag, including all intermediate labels in the
attached trees.

Choose a linear order of the vertices of $T$ extending its partial
order, and add one extra vertex. The finite ordered complex $NT$
embeds in the simplex $\Delta^V$ on this enlarged vertex set, and
every face of $NT$ is a proper subset of $V$. Both spaces are
contractible, so the monomorphism $NT\hookrightarrow\Delta^V$ is
an acyclic simplicial cofibration. Applying $c\sd^2$ gives the
Thomason trivial cofibration
\[
 D=c\sd^2NT\hookrightarrow Q=c\sd^2\Delta^V.
\]
If $P$ is fibrant, $h$ extends to a monotone $H:Q\to P$.
Inside $Q$ retain the old flags $\sigma\in D$, the flags
$\sigma^+=\sigma\cup\{V\}$, and the singleton flag $\{V\}$.
Inclusion of flags gives
$\sigma\leq\sigma^+$ and $\{V\}\leq\sigma^+$;
also $\sigma\mapsto\sigma^+$ is monotone. Thus
\[
 g(\sigma)=H(\sigma^+),\qquad c=H(\{V\})
 \quad\text{satisfy}\quad h\leq g\geq\operatorname{const}_c.
\]
By (I.9.1), $g=h$. Surjectivity of $a\mu$ implies $c\leq x$ for
every $x\in P$. Hence $c$ is a least element of $P$.

Now let $R$ be the beat-point core of an arbitrary connected finite
fibrant poset. It is a strict retract and therefore remains fibrant,
since right lifting properties are closed under retracts. It is
connected and has no upbeat points, so the preceding result gives it
a least element $c$. If $R$ had another element, a minimal element
of $R-\{c\}$ would have $c$ as greatest strict predecessor, contrary
to the absence of downbeat points. Thus $R=\{c\}$.
Fibrancy is componentwise, because each generating horn and simplex
is connected, so the result holds componentwise as stated.

The construction of $T$, all its order relations and the descending
forcing argument (I.9.1) are checked on the 22-element example of
I.1a in \texttt{code/I\_fibrancy\_checks.py}. This also produces a
finite trivial-cofibration lifting problem that the example fails.
The general argument above proves the result in every finite size;
it does not infer general fibrancy or nonfibrancy from a bounded list
of horn dimensions.
\end{proof}

\begin{claim}{I.9b}{Proved}{An exact order criterion for countable isolated identities}
Let $P$ be countable and have finite principal ideals. Enumerate its
elements $p_0,p_1,\ldots$ so that strict predecessors occur earlier.
For $p_j<y$ let $\mathcal T_{j,y}$ be the following countable tree.
A node of length $m$ is a tuple $(y_0,\ldots,y_{m-1})$ satisfying
\[
 p_i\leq y_i\quad(i<m),\qquad
 y_i\leq y_k\quad\text{if }p_i\leq p_k\text{ and }i,k<m,
\]
and $y_j=y$ when $j<m$. The tree order is extension of tuples.
Then the identity of $P$ is joined by a natural transformation in
either direction to no other endofunctor if and only if
$P$ has no downbeat points and every $\mathcal T_{j,y}$ is
well founded. This is equivalently an explicit transfinite test:
successively remove terminal nodes from each tree, at limit ordinals
intersecting the remaining node sets; every node must eventually be
removed. The test is independent of the chosen enumeration.
\end{claim}
\begin{proof}
Such an enumeration exists: starting with any listing of $P$, append
at stage $k$ the still unlisted elements of $\mathord\downarrow p_k$
in a linear order extending that finite ideal. This lists every
element and places all predecessors earlier.

If $u\leq\id_P$ moves an element, choose a moved $x$ with
$|\mathord\downarrow x|$ minimal. Every $z<x$ is fixed, so
$z=u(z)\leq u(x)<x$. Thus $u(x)$ is the greatest strict predecessor
of $x$. Conversely a downbeat deletion defines a nonidentity map
below the identity. This proves the exact downward criterion.

An infinite branch of $\mathcal T_{j,y}$ is exactly a compatible
family of the displayed tuples. Its union is a monotone map
$u:P\to P$ with $u\geq\id_P$ and $u(p_j)=y>p_j$.
Every such map conversely supplies this branch. A tree is well
founded precisely when its transfinite terminal-node deletion removes
every node: surviving nodes at stabilization each have a surviving
successor, which recursively gives a branch; a well-founded tree
admits the rank assignment
$\rho(s)=\sup\{\rho(t)+1:t\text{ is an immediate successor of }s\}$.
Since these trees are countable, all ranks used are countable
ordinals. No finite-branching compactness assertion is being used.
This proves necessity and sufficiency, and also independence of the
enumeration. Natural transformations of poset maps are exactly
pointwise inequalities. If none can leave or enter the identity,
its entire finite-zigzag component is a singleton.
\end{proof}

The finite-ideal hypothesis cannot be replaced by a claim that
absence of both kinds of beat points is sufficient. For example,
let $P=\mathbb N\times\{0,1\}$, with
$(m,i)<(n,j)$ exactly when $m<n$. Its principal ideals are finite;
neither strict lower nor strict upper sets have the required unique
extreme element. Nevertheless $(m,i)\mapsto(m+1,i)$ is a nonidentity
monotone map above the identity. On the other hand the branching
condition
\[
 x<z\quad\Longrightarrow\quad
 \text{some }y>x\text{ has no common upper bound with }z
 \tag{I.9.2}
\]
is a convenient sufficient condition for all the trees in I.9b to
be well founded. If $u\geq\id$ and $u(x)=z>x$, then $u(y)$ would
be a forbidden common upper bound. Condition (I.9.2) is sufficient,
not necessary: the six-element beat-free poset of I.8a has isolated
identity by I.1, but for $x=a_1,z=b_1$ it fails (I.9.2).

\begin{claim}{I.9c}{Proved}{Audit of INDLOCAL D.16 and D.17}
The finite-ideal and no-downbeat invariant of INDLOCAL D.16 and its
countable bifibrant example D.17 are valid. They satisfy the exact
criterion of I.9b. Their isolated identity is compatible with I.9a
because the example is infinite.
\end{claim}
\begin{proof}
We check the parts on which that conclusion depends. For
$Q=c\sd^2\Delta^n$, the horn $H=c\sd^2\Lambda_i^n$ consists
of flags containing neither $M=\{0,\ldots,n\}-\{i\}$ nor
$V=\{0,\ldots,n\}$. The greatest old subflag of a new flag
$\sigma$ is $\sigma_H=\sigma-\{M,V\}$, when nonempty.
For any attaching map $v:H\to P$, the pushout has the old objects
and separate new flags, no new-to-old arrows, and
\[
 x\leq\sigma\quad\Longleftrightarrow\quad
 \sigma_H\ne\varnothing\ \text{and }x\leq v(\sigma_H)
 \qquad(x\in P).
\]
These formulas verify the universal property even when $v$ identifies
old flags. Hence old principal ideals never change; a new ideal
is a finite collection of new subflags together with either no old
objects or the finite ideal below $v(\sigma_H)$. A new common upper
bound of two old objects gives an old common upper bound
$v(\sigma_H)$. These statements also hold for a simultaneous family,
whose new objects have separate cell labels.

A singleton new flag is minimal. If a longer new flag $\sigma$
had greatest strict predecessor $\tau$, choose a new face
$N\in\sigma\cap\{M,V\}$. The predecessor $[N]$ forces $\tau$
to be a new subflag containing $N$. Choose $S\in\sigma-\tau$.
If $S$ is new, $[S]$ is not below $\tau$. If it is old and
$|\sigma|\geq3$, the strict new subflag $[S,N]$ is not below
$\tau$. If $|\sigma|=2$, the old object $v([S])$ is below
$\sigma$ but not below the minimal new flag $[N]=\tau$.
Each case contradicts greatestness. Thus no downbeat point is
introduced. This verifies D.16 without an injectivity assumption.

Start at one point, and at each of countably many stages attach one
fresh cell for every map of every generating horn into the preceding
stage, including maps which already have fillers. There are countably
many such maps at each stage. The resulting union $F$ is countable,
and the initial inclusion is a relative generating-trivial-cofibration
cell map. Thus $F$ is cofibrant and weakly contractible. Every finite
horn map factors through a stage and is filled at the next, so $F$
is fibrant. Ideals stabilize at birth, giving finite principal ideals
and no downbeats in the union.

For $x<z$ already present, a fresh one-dimensional horn attached at
$x$ has a new maximal vertex $y>x$ whose old predecessors are
exactly $\mathord\downarrow x$. Thus $y,z$ have no common upper
bound on that pass; distinct simultaneous cells do not introduce
one. The common-upper-bound invariant preserves this absence on
every later pass and in the union. This proves (I.9.2), so I.9b
proves isolatedness of the identity. These are the essential steps
of D.16--D.17, now checked directly.
\end{proof}

\begin{claim}{I.9d}{Proved}{The core statement requires an existence hypothesis in the infinite case}
For posets with finite principal ideals, adjunction zigzags still
imply categorical homotopy equivalence. Any categorical homotopy
equivalence between posets with isolated identities is an order
isomorphism. Thus I.1 remains valid on any class of posets admitting
isolated-identity retracts reached by finite adjunction zigzags.
Such cores need not exist, even for countable posets whose principal
ideals have at most three elements. Mere absence of beat points is
also insufficient for identity rigidity, as the preceding example
shows.
\end{claim}
\begin{proof}
The unit and counit inequalities of an adjunction give the first
assertion without finiteness. For the second, if $f,g$ are homotopy
inverses, isolatedness gives $gf=\id$ and $fg=\id$, exactly as in
I.1. If the stated retract cores exist, composing with their
adjunction zigzags proves the conditional classification in both
directions and uniqueness of the cores.

For nonexistence consider the ray with vertices $0,1,2,\ldots$
and comparisons $2k<2k+1$ and $2k+2<2k+1$ only. Every connected
subposet with at least two vertices is an interval in this ray, and
its first vertex is a beat point. Hence every connected subposet
with isolated identity is a singleton. The ray is not categorically
homotopy equivalent to a point. A finite zigzag from its identity
to a constant map would, on each vertex, give a comparability path
of uniformly bounded length to the constant value, while distances
in the ray are unbounded.

Indeed it has no categorical homotopy equivalent poset with isolated
identity, whether or not initially presented as a subposet. If $R$
were such a poset, homotopy inverses $f:P\to R$, $g:R\to P$
would satisfy $fg=\id_R$ by rigidity. Thus $g$ identifies $R$
with an induced connected subposet of the ray, forcing $R$ to be
a point and yielding the contradiction just proved. Infinitely many
successive beat deletions of this ray therefore cannot be substituted
for the finite categorical homotopies in I.1.
\end{proof}
